Installing Rockbox on the \dap{} means writing Rockbox's image into the board's
flash memory and putting Rockbox's own files on the SD card.

The image is an ESP-IDF \emph{application} image, not a stand-alone firmware, so it
does not go at flash offset \fname{0}. Three pieces occupy the start of the flash:

\begin{code}
    0x000000   second-stage bootloader   (bootloader.bin)
    0x008000   partition table           (partitions.bin)
    0x010000   Rockbox                   (rockbox.esp32s3)
\end{code}

The chip's mask ROM at offset \fname{0} only copies RAM segments and jumps. It never
arms a cache MMU entry, never enables the external caches and never releases the
instruction bus; the second-stage bootloader does all of that, then maps and loads the
application. Writing Rockbox at \fname{0} therefore boots nothing -- the ROM reads
Rockbox's header as if it were a bootloader and fails.

Rockbox Utility does not offer this board yet, so the installation is done by hand. You
will need:

\begin{itemize}
  \item A Rockbox build for the MeowKit ESP32-S3 target. Building from source gives you
    all three pieces: \fname{rockbox.esp32s3} in the build directory, and
    \fname{bootloader.bin} plus \fname{partitions.bin} in the IDF application's build
    directory.
  \item A USB-C cable. The board's USB socket is wired directly to the ESP32-S3, so
    \fname{esptool} can talk to it over that connection.
  \item \fname{esptool}, version 4 or later.
  \item A FAT32-formatted microSD card. Rockbox reads its files from the card, and the
    player stores no music in internal flash.
\end{itemize}

\begin{enumerate}

\item \textbf{Connect the board, and put it in download mode.} Plug the USB-C cable
  into your computer. The serial port usually appears as \fname{/dev/ttyACM0} on Linux
  and as a \fname{COM} port on Windows.

  A board already running Rockbox or the stock firmware will not answer on that port by
  itself, because that firmware owns the USB connection. Hold \textbf{BOOT}, briefly
  press \textbf{RESET}, then release \textbf{BOOT}. That puts the chip into the ROM's
  download mode while the running firmware is still there, which is what reflashing
  needs.

\item \textbf{Check that \fname{esptool} can see the chip.}

\begin{code}
    esptool --chip esp32s3 --port /dev/ttyACM0 \fname{chip\_info}
\end{code}

  This prints the chip revision and the flash size (16~MB on this board). If it fails,
  the board is not in download mode; repeat the BOOT and RESET gesture.

\item \textbf{Write the images.} From a source build, the build's own goal is the
  simplest route: it knows the layout and the flash mode and frequency the build was
  configured for.

\begin{code}
    cd build-meowkits3
    make flash ESPPORT=/dev/ttyACM0
\end{code}

  Or write the three pieces directly, at the addresses given above:

\begin{code}
    esptool --chip esp32s3 --port /dev/ttyACM0 \fname{write\_flash} \
        0x000000 bootloader.bin \
        0x008000 partitions.bin \
        0x010000 rockbox.esp32s3
\end{code}

  \note{This replaces whatever is in those regions. If you may want the board's original
    firmware again, read the whole flash off first and keep the copy -- 16~MB covers the
    stock image and its partition table.}

\begin{code}
    esptool --chip esp32s3 --port /dev/ttyACM0 \fname{read\_flash} 0 0x1000000 stock-firmware.bin
\end{code}

  With no board attached, \fname{make espverify} runs \fname{esptool \fname{image\_info}} on the
  packed image and prints the segment table, load addresses and checksum the loader will
  see.

\item \textbf{Prepare the SD card.} Format it as FAT32, then unpack Rockbox's
  files into the root of the card. The main package contains the player's
  default theme, its plugins and codecs, and the language files:

\begin{code}
    unzip rockbox.zip -d /path/to/sd-card/
\end{code}

  This creates a \fname{.rockbox} directory at the card's root. Two further
  packages are worth unpacking into the same location:

\begin{itemize}
  \item \fname{rockbox-fonts.zip} -- the font files that many third-party
    themes need. Without them a theme that asks for a font of its own reports
    a font load error and is not displayed.
  \item The voice pack is already inside the main package: \fname{.rockbox/langs/}
    carries \fname{english.voice} and the talk clips. Nothing to install, but it is
    off until you turn it on under \fname{Settings --> Voice}, where the speaking
    rate, volume and which events are announced are set.
  \item \fname{rockbox-themes.zip} -- additional themes, if you installed any.
    Each theme ships its own backdrops and images, and the themes are selected
    in \fname{Settings --> Theme}.
\end{itemize}

\note{The \fname{.rockbox} directory is hidden on some operating systems. If
  you cannot see it in a file manager, enable viewing of hidden files before
  copying anything else onto the card, or you may end up with a second
  \fname{.rockbox} directory nested inside the first.}

\item \textbf{Or copy files over USB.} Instead of taking the card out, plug the
  board into a computer and Rockbox exposes the card as a mass-storage device. The USB
  profiles all live in \fname{Settings --> General Settings --> USB}: USB Mode
  (\setting{Mass storage} or \setting{Charge only}), USB HID on or off with its keypad
  mode, and the DAC as \setting{Never}, \setting{Always}, \setting{While charge only}
  or \setting{While mass storage}. If USB Mode is \setting{Charge only}, holding
  \emph{any} button while plugging in flips that one connection to mass storage -- the
  USB stack's insert handler inverts the charge-only decision whenever a button is
  held, and this target does not narrow the button mask that check uses, so a stick
  direction or either face button works.

\item \textbf{Start Rockbox.} Insert the card, then press and hold the
  \ButtonPower{} key to turn the board on. The player boots to the main menu.
  Rockbox reads \fname{.rockbox/config.cfg} from the card at startup, so your
  settings follow the card rather than the board.

  Lists on this player use a line padding of 4 by default. The automatic setting
  sizes each row from the panel's dot density, which on a 2.0" screen gives rows
  about 66 pixels tall -- roughly three of them visible. If you prefer larger
  rows for finger use, set it under
  \fname{Settings --> Display --> Touchscreen Settings --> Line Padding in
  Lists}; the list is a setting, so the change survives a reboot.

\end{enumerate}

To update Rockbox later, repeat step 3 with a new image. Settings and the
theme files are untouched by that step; overwriting them requires the
corresponding zip to be unpacked again.

\subsubsection{After it boots}

The panel is 320 by 240 pixels and 16-bit colour. Brightness is set under
\fname{Settings --> Display}, together with the timeout that turns the
backlight off; the panel itself follows a moment later and wakes on the next
key. There is no separate "screen off while playing" switch -- the backlight
timeout is what gets you a dark screen while audio continues.

The touchscreen covers a 3 by 3 grid over the display, and \setting{Touchscreen
Mode} decides whether a tap is reported as a point or as the button that grid
stands for. Lists size their rows from the panel's dot density, which is the
setting mentioned above; raise it if you want fewer, bigger rows.

\begin{itemize}
  \item \textbf{Sound comes out mono.} The board has one digital-to-analogue
    converter feeding a small class-D amplifier and the 3.5~mm jack; the two
    channels are mixed rather than separated. The amplifier has an enable of
    its own, exposed as the speaker setting among the sound settings, so
    audio can be routed to the jack while the speaker stays quiet.
  \item \textbf{There is no Bluetooth audio.} The chip's radio is
    Bluetooth~Low~Energy only, which has no audio transport, so headphones
    must be wired.
  \item \textbf{What it plays.} MP3, Ogg, FLAC, ALAC, Monkey's Audio,
    WavPack, TTA, Shorten, AAC, WMA and WMA Pro, ADX, AIFF, AU, WAV, Speex,
    Musepack and the chiptune formats (SID, NSF, GBS, HES, KSS, SGC, VGM,
    SPC, AY, MOD and friends) all decode on the player. WAV files up to
    192~kHz and 24-bit play through the software resampler, and DSD64 in a DSF
    container plays with its embedded album art. The processor cannot keep a
    DSD128 stream fed, so that is the ceiling.
  \item \textbf{Recording.} The board has a microphone, and \fname{Settings -->
    Recording} sets the sample rate, the input gain and whether recordings go
    to the card. Clips are written as WAV files into \fname{/recordings} on
    the SD card, and the rate can go as high as 96~kHz.
  \item \textbf{Battery.} The player reports remaining capacity from a
    coulometer -- it counts charge and discharge rather than guessing from the
    cell's voltage -- so the percentage is meaningful while charging as well as
    while playing. Charging state comes from the power management chip on the
    board, which also handles the power button: pressing and holding it turns
    the player off cleanly, which lets Rockbox finish the current track's
    buffer, save its settings and close the card before the supply is cut.
    With the screen off and nothing playing the board draws about 42~mA, so a
    dark screen is worth something.
  \item \textbf{USB.} The socket is the chip's own USB, so plugging in offers
    the profiles in \fname{Settings --> General Settings --> USB}: mass
    storage for the card, charge only, or the player as a USB audio device
    (48~kHz is the most the hardware's FIFO can carry). Holding any button
    while plugging in forces mass storage even when the mode is charge only.
\end{itemize}

Plugins need the font package described above; a plugin that draws its own
bitmaps or plays video also needs its data files on the card. Video in
particular is a codec rather than a player feature, and it asks for a small
encode -- the MPEG Player entry in the plugin documentation gives the size and
rate it wants.
